% Part: incompleteness
% Chapter: representability-in-q
% Section: undecidability

\documentclass[../../../include/open-logic-section]{subfiles}

\begin{document}

\olfileid{inc}{req}{und}
\olsection{অণির্ণেয়তা}

কোনো তত্ত্ব $\Th{T}$-কে আমরা \emph{অনির্ণেয়} বলি, যদি এমন কোনো গণনপদ্ধতি
না থাকে যা সসীমসংখ্যক ধাপের পরে এবং ব্যর্থ না হয়ে $\Th{T}$-এর ভাষার যেকোনো
বাক্য~$!A$-এর জন্য ``$\Th{T}$ কি~$!A$ প্রমাণ করে?'' প্রশ্নটির সঠিক উত্তর
দেয়। তাই পাটিগণিতের ভাষার কোনো বাক্য~$!A$ দেওয়া থাকলে
$\Th{Q} \Proves !A$ কি না নির্ণয় করে এমন গণনপদ্ধতি থাকলে এবং কেবল তখনই
$\Th{Q}$ নির্ণেয় হতো। প্রশ্নটি এভাবে আরও সুনির্দিষ্ট করা যায়:
$\Prov[\Th{Q}](y)$ সম্বন্ধটি কি পুনরাবৃত্ত, যেখানে $y$ ঠিক তখনই এই সম্বন্ধ
পূরণ করে যখন $y$ হলো $\Th{Q}$-তে প্রমাণযোগ্য কোনো বাক্যের গ্যোডেল সংখ্যা?
উত্তর: না।

\begin{thm}
$\Th{Q}$ অনির্ণেয়; অর্থাৎ সম্বন্ধ
\[
\Prov[\Th{Q}](y) \defiff \fn{Sent}(y) \land
\lexists[x][\Prf[\Th{Q}](x, y)]
\]
পুনরাবৃত্ত নয়।
\end{thm}

\begin{proof}
ধরি সেটি পুনরাবৃত্ত। তাহলে নিচের উপায়ে থামার সমস্যার সমাধান করা যেত। $e$ ও
$n$ দেওয়া থাকলে আমরা জানি, এমন একটি~$s$ আছে যাতে $T(e, n, s)$ যদি এবং কেবল
যদি $\cfind{e}(n) \fdefined$; এখানে $T$ হলো
\olref[cmp][rec][nft]{thm:kleene-nf}-এর ক্লিনির বিধেয়। $T$ আদিম পুনরাবৃত্ত
বলে $\Th{Q}$-তে কোনো সূত্র $!B_T$ সেটিকে উপস্থাপন করে; অর্থাৎ
$\Th{Q} \Proves !B_T(\num{e}, \num{n}, \num{s})$ যদি এবং কেবল যদি
$T(e, n, s)$। যদি
$\Th{Q} \Proves !B_T(\num{e}, \num{n}, \num{s})$, তবে
$\Th{Q} \Proves \lexists[y][!B_T(\num{e}, \num{n}, y)]$-ও। এমন কোনো $s$
না থাকলে প্রত্যেক~$s$-এর জন্য
$\Th{Q} \Proves \lnot !B_T(\num{e}, \num{n}, \num{s})$। কিন্তু
$\Th{Q}$ $\omega$-সঙ্গতিপূর্ণ; অর্থাৎ প্রত্যেক~$n \in \Nat$-এর জন্য
$\Th{Q} \Proves \lnot !A(\num{n})$ হলে
$\Th{Q} \Proves/ \lexists[y][!A(y)]$। এটি আমরা জানি, কারণ $\Th{Q}$-এর
স্বতঃসিদ্ধগুলি প্রমিত মডেল~$\Struct{N}$-এ সত্য। তাই
$\Th{Q} \Proves/ \lexists[y][!B_T(\num{e}, \num{n}, y)]$। অন্যভাবে বললে,
$\Th{Q} \Proves \lexists[y][!B_T(\num{e}, \num{n}, y)]$ যদি এবং কেবল যদি
এমন একটি $s$ থাকে যাতে $T(e, n, s)$; অর্থাৎ যদি এবং কেবল যদি
$\cfind{e}(n) \fdefined$। $e$ ও~$n$ থেকে
$\Gn{\lexists[y][!B_T(\num{e}, \num{n}, y)]}$ গণনা করা যায়; এই কাজটি করে
এমন আদিম পুনরাবৃত্ত অপেক্ষককে $g(e, n)$ বলি। তাই
\[
h(e, n) =
\begin{cases}
1 & \text{যদি $\Prov[\Th{Q}](g(e, n))$}\\
0 & \text{অন্যথায়}.
\end{cases}
\]
$\Prov[\Th{Q}]$ পুনরাবৃত্ত হলে এর থেকে $h$-ও পুনরাবৃত্ত হতো। কিন্তু
\olref[cmp][rec][hlt]{thm:halting-problem} অনুসারে $h$ পুনরাবৃত্ত নয়; তাই
$\Prov[\Th{Q}]$-ও পুনরাবৃত্ত হতে পারে না।
\end{proof}

\begin{cor}
প্রথম-ক্রমের যুক্তিবিদ্যা অনির্ণেয়।
\end{cor}

\begin{proof}
প্রথম-ক্রমের যুক্তিবিদ্যা নির্ণেয় হলে $\Th{Q}$-তে প্রমাণযোগ্যতাও নির্ণেয়
হতো, কারণ $\Th{Q} \Proves !A$ যদি এবং কেবল যদি $\Proves !T \lif !A$;
এখানে $!T$ হলো $\Th{Q}$-এর স্বতঃসিদ্ধগুলির সংযোজন।
\end{proof}

\end{document}
